\begin{frame}{노트북에서 확인하는 핵심 세 가지}
  \begin{enumerate}
    \item \kb{마진 경계선이 경계에 평행}하며, 그 사이가 ``마진'' —
      데이터는 마진 경계선 바깥에만 있다.
    \item \kb{서포트 벡터는 마진 경계선 위에만 있다} —
      그 밖의 점은 지워도 $w, b$가 변하지 않는다.
    \item \kb{마진 너비 $2/\|w\|$가 실제 거리와 일치} —
      경계에서 서포트 벡터까지의 유클리드 거리가 정확히
      $1/\|w\|$인지 직접 재본다.
  \end{enumerate}
  \vskip0.6em
  {\footnotesize 실습 노트북: notebooks/ml1/chapter05\_1\_svm\_margin.ipynb —
  병진(shift) 실험과 ``서포트 벡터 아닌 점 이동'' 실험까지 포함.}
\end{frame}
